%% Generated by Verifpal 1.3.7. Compile with: tectonic report.tex
\documentclass[11pt]{article}
\usepackage[margin=2.2cm]{geometry}
%% --- BEGIN verifpal preamble ---
% SPDX-FileCopyrightText: (c) 2019-2026 Nadim Kobeissi <nadim@symbolic.software>
% SPDX-License-Identifier: GPL-3.0-only
\usepackage{iftex}
\ifPDFTeX
  \usepackage[T1]{fontenc}
  \usepackage[utf8]{inputenc}
\fi
\usepackage{amsmath}
\usepackage{amssymb}
\usepackage{xcolor}
\usepackage{tikz}
\usetikzlibrary{arrows.meta,backgrounds,calc}
\usepackage{booktabs}
\usepackage{longtable}
\usepackage{enumitem}
\usepackage{caption}
\usepackage{listings}
\usepackage{microtype}
\usepackage[hidelinks]{hyperref}

\definecolor{vpink}{RGB}{26,30,36}
\definecolor{vpslate}{RGB}{92,104,120}
\definecolor{vprule}{RGB}{196,203,212}
\definecolor{vpaccent}{RGB}{158,32,44}
\definecolor{vpband}{RGB}{243,245,248}
\definecolor{vptint}{RGB}{252,243,243}

\captionsetup{font=small,labelfont=bf,skip=7pt}

% Terms. Redefine any of these to restyle every figure and every trace at once.
\newcommand{\vpus}{\kern0.06em\rule[-0.22ex]{0.34em}{0.045em}\kern0.06em}
\newcommand{\vpprim}[2]{\mathsf{#1}_{\mathsf{#2}}}
\newcommand{\vpprimo}[1]{\mathsf{#1}}
\newcommand{\vpconst}[1]{\mathit{#1}}
\newcommand{\vpsession}[2]{#1^{\##2}}
\newcommand{\vpscenario}[2]{#1^{@#2}}
\newcommand{\vpguard}[1]{[#1]}
\newcommand{\vpcaps}[1]{^{[\mathrm{#1}]}}
\newcommand{\vpchecked}{^{?}}
\newcommand{\vpsep}{\allowbreak\,}
\newcommand{\vpterm}[1]{\ensuremath{#1}}
\newcommand{\vpliteral}[1]{\texttt{#1}}
\newcommand{\vpactor}[1]{\textsf{#1}}
\newcommand{\vpdagger}[1]{\ensuremath{{#1}^{\dagger}}}
\newcommand{\vpguarded}[1]{[#1]}
\newcommand{\vpverdictpass}{\textcolor{vpslate}{\textsf{holds}}}
\newcommand{\vpverdictfail}{\textcolor{vpaccent}{\textsf{attack}}}

% Diagrams. Three colours only: ink, grey, and one accent. The accent is spent
% on what is security-relevant -- the attacker's lane, forged traffic, and the
% checked primitives at which a principal can halt -- and on nothing else.
\newsavebox{\vpbox}
\newsavebox{\vpmeasure}
\newdimen\vptmp
\newcounter{vplanes}
\newif\ifvpfirstline
\newlength{\vplanesep}
\newlength{\vpboxw}
\newlength{\vpnumgut}\setlength{\vpnumgut}{0.62cm}
\def\vprowsep{0.80}

\tikzset{
  vpactorplate/.style={draw=vprule,line width=0.6pt,fill=vpband,rounded corners=2pt,
    minimum width=2.1cm,minimum height=0.52cm,font=\small\bfseries,text=vpink,inner sep=2pt},
  vpactorplateadv/.style={vpactorplate,draw=vpaccent!50,fill=white,text=vpaccent},
  vplife/.style={draw=vprule,line width=0.6pt},
  vplifeadv/.style={draw=vpaccent!35,line width=0.6pt},
  vparrow/.style={-{Stealth[length=1.9mm,width=1.5mm]},line width=0.55pt},
  vpplain/.style={draw=vpink!75},
  vpforged/.style={draw=vpaccent,line width=0.7pt},
  vpreplay/.style={draw=vpaccent,line width=0.7pt,dash pattern=on 3pt off 1.6pt},
  vpmuted/.style={draw=vprule,line width=0.55pt},
  vplabel/.style={font=\footnotesize,fill=white,inner sep=2.5pt,text=vpink},
  vplabeladv/.style={vplabel,text=vpaccent},
  vpactbox/.style={draw=vprule,fill=white,line width=0.5pt,rounded corners=1.5pt,
    inner xsep=5pt,inner ysep=4pt,font=\scriptsize,text=vpink!85,align=left},
  vpleakbox/.style={vpactbox,draw=vpslate!45,fill=vpband},
  vpnote/.style={draw=vprule,fill=vpband,line width=0.5pt,rounded corners=1.5pt,
    font=\scriptsize,text=vpslate,inner sep=3pt},
  vpnoteadv/.style={vpnote,draw=vpaccent!45,fill=white,text=vpaccent},
  vpnum/.style={font=\scriptsize,text=vpslate,inner sep=0pt,anchor=east},
  vpnumadv/.style={vpnum,text=vpaccent},
  vpnumwide/.style={vpnum},
  vpnumwideadv/.style={vpnumadv},
  vpphaselabel/.style={font=\scriptsize\scshape,text=vpslate,inner sep=2pt},
}

% Rows advance a running y, so a tall box pushes what follows down rather than
% overlapping it.
\newcommand{\vpadvance}[1]{%
  \pgfmathparse{\vpyc-(#1)/2}\global\edef\vpy{\pgfmathresult}%
  \pgfmathparse{\vpyc-(#1)}\global\edef\vpyc{\pgfmathresult}}

\newcommand{\vplane}[2]{\node[vpactorplate] at (#1,0) {#2};}
\newcommand{\vplaneadv}[2]{%
  \gdef\vpadvlane{#1}\node[vpactorplateadv] at (#1,0) {#2};}

% Step numbers sit in the left margin, clear of every lane.
\newcommand{\vpstep}[2]{%
  \if\relax\detokenize{#2}\relax\else
    \node[#1] at (0,\vpy) [xshift=-0.5\vpboxw-\vpnumgut] {#2};
  \fi}

\newcommand{\vphop}[7]{%
  \vpadvance{\vprowsep}%
  \ifnum#2>-1\relax
    \draw[vpmuted] (#1,\vpy) -- (#2,\vpy);
    \draw[vparrow,#4] (#2,\vpy) -- (#3,\vpy) node[vplabel,midway]{#7};
  \else
    \draw[vparrow,#4] (#1,\vpy) -- (#3,\vpy) node[vplabel,midway]{#7};
  \fi
  \vpstep{#5}{#6}}

\newcommand{\vphopforged}[7]{%
  \vpadvance{\vprowsep}%
  \ifnum#2>-1\relax
    \draw[vpmuted] (#1,\vpy) -- (#2,\vpy);
    \draw[vparrow,#4] (#2,\vpy) -- (#3,\vpy) node[vplabeladv,midway]{#7};
  \else
    \draw[vparrow,#4] (#1,\vpy) -- (#3,\vpy) node[vplabeladv,midway]{#7};
  \fi
  \vpstep{#5}{#6}}

\newcommand{\vpphase}[1]{%
  \vpadvance{0.62}%
  \begin{scope}[on background layer]
    \draw[vprule,line width=0.5pt,dash pattern=on 2.4pt off 2.4pt]
      ($(0,\vpy)+(-0.5\vpboxw-\vpnumgut,0)$) -- ($({\thevplanes-1},\vpy)+(0.5\vpboxw,0)$);
  \end{scope}
  \node[vpphaselabel,anchor=west,fill=white,inner xsep=4pt]
    at (0,\vpy) [xshift=-0.5\vpboxw-\vpnumgut] {phase #1};}

% Every box is the same width, so the rows form aligned columns. Its height
% comes from the line count rather than from measuring: pgfpicture selects
% \nullfont, so a box measured inside one comes back the wrong size, and a
% line count also keeps the vertical rhythm regular.
\def\vplineh{0.335}
\newcommand{\vpboxrow}[4]{%
  \pgfmathparse{#1*\vplineh+0.34}\edef\vpboxh{\pgfmathresult}%
  \vpadvance{\vpboxh}%
  \global\vpfirstlinetrue
  \node[#2,text width=\vpboxw] at (#3,\vpy) {#4};}

\newcommand{\vpactline}[1]{\ifvpfirstline\global\vpfirstlinefalse\else\\\fi#1}
\newcommand{\vpgen}[1]{%
  \vpactline{{\color{vpslate}\scshape new}\ #1}}
\newcommand{\vpcompute}[2]{%
  \vpactline{#1\ {\color{vprule}\ensuremath{\leftarrow}}\ #2}}
\newcommand{\vpcomputechecked}[2]{%
  \vpactline{#1\ {\color{vprule}\ensuremath{\leftarrow}}\ {\color{vpaccent}#2}}}
\newcommand{\vpalias}[1]{\vpactline{#1}}
\newcommand{\vpcheckonly}[1]{\vpactline{{\color{vpaccent}#1}}}

\newcommand{\vpactivity}[3]{\vpboxrow{#2}{vpactbox}{#1}{#3}}

\newcommand{\vpleak}[3]{\vpboxrow{#2}{vpleakbox}{#1}{#3}}

\newcommand{\vpmark}[5]{%
  \vpadvance{\vprowsep}%
  \node[#2] at (#1,\vpy) {#5};%
  \vpstep{#3}{#4}}

\newcommand{\vprun}[4]{%
  \vpadvance{\vprowsep}%
  \node[vpnote] at (#1,\vpy) {#4};%
  \vpstep{#2}{#3}}

\newenvironment{vpdiagram}[1]{%
  \setcounter{vplanes}{#1}%
  \gdef\vpadvlane{-1}%
  \gdef\vpyc{-0.46}%
  \pgfmathsetlength{\vplanesep}{max(4.6cm,(\linewidth-1.4cm)/max(1,#1-1))}%
  \pgfmathsetlength{\vpboxw}%
    {max(3.2cm,min((\linewidth-1.2cm)/#1-0.5cm,\vplanesep-0.9cm))}%
  \begin{lrbox}{\vpbox}%
  \begin{tikzpicture}[x=\vplanesep,y=1cm]%
}{%
  \begin{scope}[on background layer]
    \ifnum\vpadvlane>-1\relax
      \draw[vplifeadv] (\vpadvlane,-0.28) -- (\vpadvlane,{\vpyc-0.1});
    \fi
    \foreach \vpi in {0,...,\the\numexpr\thevplanes-1\relax} {
      \ifnum\vpi=\vpadvlane\relax\else
        \draw[vplife] (\vpi,-0.28) -- (\vpi,{\vpyc-0.1});
      \fi
    }
  \end{scope}
  \end{tikzpicture}%
  \end{lrbox}%
  \pgfmathsetmacro{\vpscale}{min(1,\linewidth/\wd\vpbox,
    0.93*\textheight/(\ht\vpbox+\dp\vpbox))}%
  \scalebox{\vpscale}{\usebox{\vpbox}}}

% Traces.
\newlist{vptrace}{enumerate}{1}
\setlist[vptrace]{label=\protect\vptracenum{\arabic*},leftmargin=2.2em,
  itemsep=1.5pt,topsep=5pt,parsep=0pt,font=\footnotesize,before=\footnotesize}
\newcommand{\vptracenum}[1]{%
  \tikz[baseline=(vpn.base)]{\node[circle,fill=vpslate,text=white,font=\tiny,
    inner sep=1.3pt] (vpn) {#1};}}
% Flush right, breaking to its own line when the item fills the measure.
\newcommand{\vpon}[2]{%
  \unskip\nobreak\hfil\penalty50\hskip1em\null\nobreak\hfil
  {\scriptsize\textsf{#1}\,\ensuremath{\rightarrow}\,\textsf{#2}}%
  \begingroup\parfillskip=0pt\par\endgroup}

% Model source listings.
\lstdefinelanguage{Verifpal}{
  morekeywords={attacker,active,passive,principal,knows,generates,leaks,phase,
    queries,scenarios,precondition,public,private,password,confidentiality,
    authentication,freshness,unlinkability,equivalence,weak,forgeable,malleable,from},
  morekeywords=[2]{ASSERT,CONCAT,SPLIT,HASH,PW_HASH,HKDF,AEAD_ENC,AEAD_DEC,ENC,DEC,
    MAC,PUBKEY,DH_KEX,SIGN,SIGNVERIF,PKE_ENC,PKE_DEC,SHAMIR_SPLIT,SHAMIR_JOIN,
    RINGSIGN,RINGSIGNVERIF,BLIND,UNBLIND,KEM_ENCAP,KEM_DECAP},
  sensitive=false,
  morecomment=[l]{//},
  morecomment=[s]{/*}{*/},
  literate={->}{{$\rightarrow$\,}}3 {→}{{$\rightarrow$\,}}2,
}
\lstset{
  language=Verifpal,
  basicstyle=\ttfamily\footnotesize,
  keywordstyle=\color{vpaccent},
  keywordstyle=[2]\color{vpslate}\bfseries,
  commentstyle=\color{vpslate}\itshape,
  showstringspaces=false,
  breaklines=true,
  columns=fullflexible,
  frame=leftline,
  framesep=6pt,
  rulecolor=\color{vprule},
  xleftmargin=10pt,
  numbers=left,
  numberstyle=\tiny\color{vprule},
  numbersep=8pt,
}
%% --- END verifpal preamble ---

\title{Verifpal analysis of 01-signature-bypass.vp}
\author{}
\date{\today}

\begin{document}
\maketitle

\begin{abstract}
\noindent This report describes the analysis of the Verifpal model \texttt{01-signature-bypass.vp} against an active attacker, with each principal running 2 concurrent sessions. Of its 2 queries, 2 were broken. Every attack reported here is a witness against the model as written. A query reported as holding means that this search found no attack at these parameters, which is not a proof that none exists.
\end{abstract}



\section{\texorpdfstring{\texttt{01-signature-bypass.vp}}{01-signature-bypass.vp}}
\label{sec:01-signature-bypass-vp-0}


\noindent Analysed against an active attacker at 2 sessions per principal, in 170 ms. The result code is \texttt{a1c1}.


\subsection{Protocol}
%% --- BEGIN verifpal figure: 01-signature-bypass-vp-0-protocol ---
%% Self-contained: needs only the verifpal preamble block.
\begin{figure}[htbp]
\centering
\begin{vpdiagram}{2}
\vplane{0}{Wallet}
\vplane{1}{Dapp}
\vpactivity{0}{1}{\vpcompute{\vpterm{\vpconst{gw}}}{\vpterm{\vpprimo{PUBKEY}(\vpconst{wsk})}}}
\vphop{0}{-1}{1}{vpplain}{vpnum}{1}{\vpguarded{\vpterm{\vpconst{gw}}}}
\vpactivity{1}{4}{\vpgen{\vpterm{\vpconst{x}}, \vpterm{\vpconst{nonce}}, \vpterm{\vpconst{txdata}}}\vpcompute{\vpterm{\vpconst{gx}}}{\vpterm{\vpprimo{PUBKEY}(\vpconst{x})}}\vpcompute{\vpterm{\vpconst{req}}}{\vpterm{\vpprim{AEAD}{ENC}(\vpconst{symkey},\vpsep \vpprimo{CONCAT}(\vpconst{c\vpus auth},\vpsep \vpconst{gx},\vpsep \vpconst{nonce},\vpsep \vpconst{c\vpus dom}),\vpsep \vpconst{c\vpus ad})}}}
\vpleak{1}{1}{Dapp leaks symkey}
\vphop{1}{-1}{0}{vpplain}{vpnum}{2}{\vpterm{\vpconst{req}}}
\vpactivity{0}{9}{\vpgen{\vpterm{\vpconst{y}}}\vpcompute{\vpterm{\vpconst{gy}}}{\vpterm{\vpprimo{PUBKEY}(\vpconst{y})}}\vpcomputechecked{\vpterm{\vpconst{d\vpus w}}}{\vpterm{\vpprim{AEAD}{DEC}\vpchecked{}(\vpconst{symkey},\vpsep \vpconst{req},\vpsep \vpconst{c\vpus ad})}}\vpcomputechecked{\vpterm{\vpconst{tag\vpus w}}, \vpterm{\vpconst{gx\vpus w}}, \vpterm{\vpconst{nonce\vpus w}}, \vpterm{\vpconst{dom\vpus w}}}{\vpterm{\vpprimo{SPLIT}\vpchecked{}(\vpconst{d\vpus w})}}\vpcheckonly{\vpterm{\vpprimo{ASSERT}\vpchecked{}(\vpconst{tag\vpus w},\vpsep \vpconst{c\vpus auth})}}\vpcompute{\vpterm{\vpconst{cp}}}{\vpterm{\vpprimo{CONCAT}(\vpconst{gw},\vpsep \vpconst{dom\vpus w},\vpsep \vpconst{nonce\vpus w},\vpsep \vpconst{c\vpus res})}}\vpcompute{\vpterm{\vpconst{cs}}}{\vpterm{\vpprimo{SIGN}(\vpconst{wsk},\vpsep \vpconst{cp})}}\vpcompute{\vpterm{\vpconst{r}}}{\vpterm{\vpprimo{HKDF}(\vpconst{nil},\vpsep \vpprim{DH}{KEX}(\vpconst{gx\vpus w},\vpsep \vpconst{y}),\vpsep \vpconst{nil})}}\vpcompute{\vpterm{\vpconst{resp}}}{\vpterm{\vpprim{AEAD}{ENC}(\vpconst{r},\vpsep \vpprimo{CONCAT}(\vpconst{cp},\vpsep \vpconst{cs},\vpsep \vpconst{gy}),\vpsep \vpconst{c\vpus ad})}}}
\vphop{0}{-1}{1}{vpplain}{vpnum}{3}{\vpdagger{\vpterm{\vpconst{gy}}},\, \vpdagger{\vpterm{\vpconst{resp}}}}
\vpactivity{1}{7}{\vpcompute{\vpterm{\vpconst{r\vpus d}}}{\vpterm{\vpprimo{HKDF}(\vpconst{nil},\vpsep \vpprim{DH}{KEX}(\vpconst{gy},\vpsep \vpconst{x}),\vpsep \vpconst{nil})}}\vpcomputechecked{\vpterm{\vpconst{d\vpus d}}}{\vpterm{\vpprim{AEAD}{DEC}\vpchecked{}(\vpconst{r\vpus d},\vpsep \vpconst{resp},\vpsep \vpconst{c\vpus ad})}}\vpcomputechecked{\vpterm{\vpconst{cp\vpus d}}, \vpterm{\vpconst{cs\vpus d}}, \vpterm{\vpconst{gy\vpus pay}}}{\vpterm{\vpprimo{SPLIT}\vpchecked{}(\vpconst{d\vpus d})}}\vpcomputechecked{\vpterm{\vpconst{iss\vpus d}}, \vpterm{\vpconst{dom\vpus d}}, \vpterm{\vpconst{non\vpus d}}, \vpterm{\vpconst{res\vpus d}}}{\vpterm{\vpprimo{SPLIT}\vpchecked{}(\vpconst{cp\vpus d})}}\vpcheckonly{\vpterm{\vpprimo{ASSERT}\vpchecked{}(\vpconst{iss\vpus d},\vpsep \vpconst{gw})}}\vpcompute{\vpterm{\vpconst{sess\vpus d}}}{\vpterm{\vpprimo{HKDF}(\vpconst{nil},\vpsep \vpprim{DH}{KEX}(\vpconst{gy\vpus pay},\vpsep \vpconst{x}),\vpsep \vpconst{nil})}}\vpcompute{\vpterm{\vpconst{sessreq}}}{\vpterm{\vpprim{AEAD}{ENC}(\vpconst{sess\vpus d},\vpsep \vpconst{txdata},\vpsep \vpconst{c\vpus ad})}}}
\vphop{1}{-1}{0}{vpplain}{vpnum}{4}{\vpterm{\vpconst{sessreq}}}
\vpactivity{0}{1}{\vpcheckonly{\vpterm{\vpprim{AEAD}{DEC}\vpchecked{}(\vpconst{r},\vpsep \vpconst{sessreq},\vpsep \vpconst{c\vpus ad})}}}
\end{vpdiagram}
\caption{The protocol as written. Guarded values are bracketed; a dagger marks every value an attack below substitutes or replays.}
\label{fig:01-signature-bypass-vp-0-protocol}
\end{figure}
%% --- END verifpal figure: 01-signature-bypass-vp-0-protocol ---


\subsection{Verdicts}
%% --- BEGIN verifpal table: 01-signature-bypass-vp-0-verdicts ---
\begin{center}
\captionof{table}{Every query in the model, with the verdict this run reached.}
\label{tab:01-signature-bypass-vp-0-verdicts}
\begin{tabular}{@{}llll@{}}
\toprule
Query & Verdict & Scope of the verdict & Line \\
\midrule
\textsf{authentication?}~\vpactor{Wallet} \ensuremath{\rightarrow} \vpactor{Dapp}: \vpterm{\vpconst{resp}} & \vpverdictfail & a witness, shown below & 182 \\\textsf{confidentiality?}~\vpterm{\vpconst{txdata}} & \vpverdictfail & a witness, shown below & 184 \\\bottomrule
\end{tabular}
\end{center}
%% --- END verifpal table: 01-signature-bypass-vp-0-verdicts ---


\subsection{Attacks}

\subsubsection{\texorpdfstring{\textsf{authentication?}~\vpactor{Wallet} \ensuremath{\rightarrow} \vpactor{Dapp}: \vpterm{\vpconst{resp}}}{authentication? Wallet -\textgreater{} Dapp: resp}}
\noindent \vpterm{\vpconst{resp}} (\vpterm{\vpprim{AEAD}{ENC}(\vpconst{r\vpus d},\vpsep \vpconst{d\vpus d},\vpsep \vpconst{c\vpus ad})}), sent by \vpactor{Attacker} and not by \vpactor{Wallet}, is successfully used in \vpterm{\vpprim{AEAD}{DEC}\vpchecked{}(\vpconst{r\vpus d},\vpsep \vpconst{resp},\vpsep \vpconst{c\vpus ad})} within \vpactor{Dapp}'s state.

%% --- BEGIN verifpal figure: 01-signature-bypass-vp-0-attack-0 ---
%% Self-contained: needs only the verifpal preamble block.
\begin{center}
\begin{vpdiagram}{3}
\vplane{0}{Dapp}
\vplane{1}{Wallet}
\vplaneadv{2}{Attacker}
\vpleak{0}{1}{Dapp leaks symkey}
\vprun{2}{vpnumwide}{1--13}{computes}
\vphopforged{1}{2}{0}{vpforged}{vpnumadv}{14}{\vpdagger{\vpterm{\vpconst{gy}}},\, \vpdagger{\vpterm{\vpconst{resp}}}}
\vpmark{0}{vpnote}{vpnumadv}{15}{check passes}
\vpmark{0}{vpnote}{vpnumadv}{16}{check passes}
\vpmark{0}{vpnote}{vpnumadv}{17}{check passes}
\vpmark{0}{vpnote}{vpnumadv}{18}{check passes}
\end{vpdiagram}
\captionof{figure}{How the attacker breaks \textnormal{\textsf{authentication?}~\vpactor{Wallet} \ensuremath{\rightarrow} \vpactor{Dapp}: \vpterm{\vpconst{resp}}}. Substituted values carry a dagger.}
\label{fig:01-signature-bypass-vp-0-attack-0}
\end{center}
%% --- END verifpal figure: 01-signature-bypass-vp-0-attack-0 ---


\begin{vptrace}
\item During \vpactor{Wallet}'s run (phase 0), attacker constructs \vpterm{\vpprimo{PUBKEY}(\vpconst{nil})}.
\item During \vpactor{Wallet}'s run (phase 0), attacker observes \vpterm{\vpconst{req}} on the wire.
\item During \vpactor{Wallet}'s run (phase 0), attacker is handed \vpterm{\vpconst{symkey}} by a leaks declaration.
\item During \vpactor{Wallet}'s run (phase 0), attacker opens \vpterm{\vpconst{req}} with \vpterm{\vpconst{symkey}}, obtaining \vpterm{\vpconst{d\vpus w}}.
\item During \vpactor{Wallet}'s run (phase 0), attacker splits \vpterm{\vpconst{d\vpus w}} and takes \vpterm{\vpconst{gx}}.
\item \vpactor{Attacker} constructs \vpterm{\vpprim{DH}{KEX}(\vpconst{gx},\vpsep \vpconst{nil})}.
\item \vpactor{Attacker} constructs \vpterm{\vpprimo{HKDF}(\vpconst{nil},\vpsep \vpprim{DH}{KEX}(\vpconst{gx},\vpsep \vpconst{nil}),\vpsep \vpconst{nil})}.
\item During \vpactor{Wallet}'s run (phase 0), attacker observes \vpterm{\vpprim{AEAD}{ENC}(\vpprimo{HKDF}(\vpconst{nil},\vpsep \vpprim{DH}{KEX}(\vpconst{gy\vpus pay},\vpsep \vpconst{nil}),\vpsep \vpconst{nil}),\vpsep \vpconst{d\vpus d},\vpsep \vpconst{c\vpus ad})} on the wire.
\item During \vpactor{Wallet}'s run (phase 0), attacker observes \vpterm{\vpprimo{PUBKEY}(\vpconst{y})} on the wire.
\item During \vpactor{Wallet}'s run (phase 0), attacker constructs \vpterm{\vpprim{DH}{KEX}(\vpconst{gy\vpus pay},\vpsep \vpconst{nil})}.
\item During \vpactor{Wallet}'s run (phase 0), attacker constructs \vpterm{\vpprimo{HKDF}(\vpconst{nil},\vpsep \vpprim{DH}{KEX}(\vpconst{gy\vpus pay},\vpsep \vpconst{nil}),\vpsep \vpconst{nil})}.
\item During \vpactor{Wallet}'s run (phase 0), attacker opens \vpterm{\vpprim{AEAD}{ENC}(\vpprimo{HKDF}(\vpconst{nil},\vpsep \vpprim{DH}{KEX}(\vpconst{gy\vpus pay},\vpsep \vpconst{nil}),\vpsep \vpconst{nil}),\vpsep \vpconst{d\vpus d},\vpsep \vpconst{c\vpus ad})} with \vpterm{\vpprimo{HKDF}(\vpconst{nil},\vpsep \vpprim{DH}{KEX}(\vpconst{gy\vpus pay},\vpsep \vpconst{nil}),\vpsep \vpconst{nil})}, obtaining \vpterm{\vpprimo{CONCAT}(\vpconst{cp},\vpsep \vpconst{cs},\vpsep \vpconst{gy\vpus pay})}.
\item \vpactor{Attacker} constructs \vpterm{\vpprim{AEAD}{ENC}(\vpconst{r\vpus d},\vpsep \vpconst{d\vpus d},\vpsep \vpconst{c\vpus ad})}.
\item \vpactor{Attacker} replaces \vpterm{\vpconst{gy}}, \vpterm{\vpconst{resp}} (sent by \vpactor{Wallet} to \vpactor{Dapp}) with \vpterm{\vpprimo{PUBKEY}(\vpconst{nil})}, \vpterm{\vpprim{AEAD}{ENC}(\vpconst{r\vpus d},\vpsep \vpconst{d\vpus d},\vpsep \vpconst{c\vpus ad})}. (\vpterm{\vpconst{gy}} was \vpterm{\vpprimo{PUBKEY}(\vpconst{y})}; \vpterm{\vpconst{resp}} was \vpterm{\vpprim{AEAD}{ENC}(\vpconst{r},\vpsep \vpprimo{CONCAT}(\vpconst{cp},\vpsep \vpconst{cs},\vpsep \vpconst{gy}),\vpsep \vpconst{c\vpus ad})})\vpon{Wallet}{Dapp}
\item \vpactor{Dapp}'s \vpterm{\vpprim{AEAD}{DEC}\vpchecked{}(\vpconst{r\vpus d},\vpsep \vpprim{AEAD}{ENC}(\vpconst{r\vpus d},\vpsep \vpprimo{CONCAT}(\vpconst{cp},\vpsep \vpconst{cs},\vpsep \vpconst{gy\vpus pay}),\vpsep \vpconst{c\vpus ad}),\vpsep \vpconst{c\vpus ad})} passes --- the attacker controls one of its inputs.
\item \vpactor{Dapp}'s \vpterm{\vpprimo{SPLIT}\vpchecked{}(\vpconst{d\vpus d})} passes --- the attacker controls one of its inputs.
\item \vpactor{Dapp}'s \vpterm{\vpprimo{SPLIT}\vpchecked{}(\vpconst{cp\vpus d})} passes --- the attacker controls one of its inputs.
\item \vpactor{Dapp}'s \vpterm{\vpprimo{ASSERT}\vpchecked{}(\vpconst{iss\vpus d},\vpsep \vpconst{gw})} passes --- the attacker controls one of its inputs.
\end{vptrace}

\subsubsection{\texorpdfstring{\textsf{confidentiality?}~\vpterm{\vpconst{txdata}}}{confidentiality? txdata}}
\noindent \vpterm{\vpconst{txdata}} (\vpterm{\vpconst{txdata}}) is obtained by \vpactor{Attacker}.

%% --- BEGIN verifpal figure: 01-signature-bypass-vp-0-attack-1 ---
%% Self-contained: needs only the verifpal preamble block.
\begin{center}
\begin{vpdiagram}{3}
\vplane{0}{Dapp}
\vplane{1}{Wallet}
\vplaneadv{2}{Attacker}
\vpleak{0}{1}{Dapp leaks symkey}
\vprun{2}{vpnumwide}{1--12}{computes}
\vphopforged{1}{2}{0}{vpforged}{vpnumadv}{13}{\vpdagger{\vpterm{\vpconst{gy}}},\, \vpdagger{\vpterm{\vpconst{resp}}}}
\vpmark{0}{vpnote}{vpnumadv}{14}{check passes}
\vpmark{0}{vpnote}{vpnumadv}{15}{check passes}
\vpmark{0}{vpnote}{vpnumadv}{16}{check passes}
\vpmark{0}{vpnote}{vpnumadv}{17}{check passes}
\vprun{2}{vpnumwide}{18--19}{computes}
\end{vpdiagram}
\captionof{figure}{How the attacker breaks \textnormal{\textsf{confidentiality?}~\vpterm{\vpconst{txdata}}}. Substituted values carry a dagger.}
\label{fig:01-signature-bypass-vp-0-attack-1}
\end{center}
%% --- END verifpal figure: 01-signature-bypass-vp-0-attack-1 ---


\begin{vptrace}
\item \vpactor{Attacker} constructs \vpterm{\vpprimo{PUBKEY}(\vpconst{nil})}.
\item \vpactor{Attacker} observes \vpterm{\vpconst{req}} on the wire.
\item \vpactor{Attacker} is handed \vpterm{\vpconst{symkey}} by a leaks declaration.
\item \vpactor{Attacker} opens \vpterm{\vpconst{req}} with \vpterm{\vpconst{symkey}}, obtaining \vpterm{\vpconst{d\vpus w}}.
\item \vpactor{Attacker} splits \vpterm{\vpconst{d\vpus w}} and takes \vpterm{\vpconst{gx}}.
\item \vpactor{Attacker} constructs \vpterm{\vpprim{DH}{KEX}(\vpconst{gx},\vpsep \vpconst{nil})}.
\item \vpactor{Attacker} constructs \vpterm{\vpprimo{HKDF}(\vpconst{nil},\vpsep \vpprim{DH}{KEX}(\vpconst{gx},\vpsep \vpconst{nil}),\vpsep \vpconst{nil})}.
\item \vpactor{Attacker} observes \vpterm{\vpconst{gw}} on the wire.
\item \vpactor{Attacker} splits \vpterm{\vpconst{d\vpus w}} and takes \vpterm{\vpconst{nonce}}.
\item \vpactor{Attacker} constructs \vpterm{\vpconst{cp}}.
\item \vpactor{Attacker} constructs \vpterm{\vpprimo{CONCAT}(\vpconst{cp},\vpsep \vpconst{nil},\vpsep \vpprimo{PUBKEY}(\vpconst{nil}))}.
\item \vpactor{Attacker} constructs \vpterm{\vpprim{AEAD}{ENC}(\vpconst{r\vpus d},\vpsep \vpconst{d\vpus d},\vpsep \vpconst{c\vpus ad})}.
\item \vpactor{Attacker} replaces \vpterm{\vpconst{gy}}, \vpterm{\vpconst{resp}} (sent by \vpactor{Wallet} to \vpactor{Dapp}) with \vpterm{\vpprimo{PUBKEY}(\vpconst{nil})}, \vpterm{\vpprim{AEAD}{ENC}(\vpconst{r\vpus d},\vpsep \vpconst{d\vpus d},\vpsep \vpconst{c\vpus ad})}. (\vpterm{\vpconst{gy}} was \vpterm{\vpprimo{PUBKEY}(\vpconst{y})}; \vpterm{\vpconst{resp}} was \vpterm{\vpprim{AEAD}{ENC}(\vpconst{r},\vpsep \vpprimo{CONCAT}(\vpconst{cp},\vpsep \vpconst{cs},\vpsep \vpconst{gy}),\vpsep \vpconst{c\vpus ad})})\vpon{Wallet}{Dapp}
\item \vpactor{Dapp}'s \vpterm{\vpprim{AEAD}{DEC}\vpchecked{}(\vpconst{r\vpus d},\vpsep \vpprim{AEAD}{ENC}(\vpconst{r\vpus d},\vpsep \vpprimo{CONCAT}(\vpconst{cp},\vpsep \vpconst{nil},\vpsep \vpprimo{PUBKEY}(\vpconst{nil})),\vpsep \vpconst{c\vpus ad}),\vpsep \vpconst{c\vpus ad})} passes --- the attacker controls one of its inputs.
\item \vpactor{Dapp}'s \vpterm{\vpprimo{SPLIT}\vpchecked{}(\vpconst{d\vpus d})} passes --- the attacker controls one of its inputs.
\item \vpactor{Dapp}'s \vpterm{\vpprimo{SPLIT}\vpchecked{}(\vpconst{cp\vpus d})} passes --- the attacker controls one of its inputs.
\item \vpactor{Dapp}'s \vpterm{\vpprimo{ASSERT}\vpchecked{}(\vpconst{iss\vpus d},\vpsep \vpconst{gw})} passes --- the attacker controls one of its inputs.
\item \vpactor{Attacker} observes \vpterm{\vpconst{sessreq}} on the wire.
\item \vpactor{Attacker} opens \vpterm{\vpconst{sessreq}} with \vpterm{\vpconst{r\vpus d}}, obtaining \vpterm{\vpconst{txdata}}.
\end{vptrace}



\subsection{Scope and assumptions}
\noindent Every verdict above was reached against an active attacker, with each principal running 2 concurrent sessions, over exactly the model as written. An attack is a witness and stands on its own. A query reported as holding says only that this search found no attack at those parameters: the search space this engine defines was explored, which is never the space of all attacks.




\appendix
\section{Model source}

\subsection{\texorpdfstring{\texttt{01-signature-bypass.vp}}{01-signature-bypass.vp}}
\label{lst:01-signature-bypass-vp-0}
%% --- BEGIN verifpal listing: 01-signature-bypass-vp-0 ---
\begin{lstlisting}
// FINDING 1 — a CACAO whose signature does not verify still establishes a
// persisted, subscribed session, for an account the attacker names.
//
// Expected: a1c1 at one session and at two.
//
// ---------------------------------------------------------------------------
// SOURCE OF EVERY MODELLING DECISION BELOW
// walletconnect-monorepo @ master
// ---------------------------------------------------------------------------
//
// utils/src/crypto.ts
//   deriveSymKey(priv, pub) = hkdf(sha256, x25519(priv, pub), undefined,
//                                  undefined, 32)          -- no salt, no info
//   hashKey(k)              = sha256(k)
//   type 1 envelope         = type | senderPublicKey | iv | sealed
//   encrypt()               = chacha20poly1305(symKey, iv).encrypt(msg)
//                             with NO associated data, so senderPublicKey is
//                             cleartext and unauthenticated.
//
// core/src/controllers/crypto.ts
//   encode(): for a type 1 envelope the symKey is NOT the topic's key --
//     topic = await this.generateSharedKey(senderPublicKey, receiverPublicKey)
//   so the response is sealed under the very key that becomes the session key.
//   decode(): the receiver derives that key from the senderPublicKey carried
//     IN the envelope:
//       generateSharedKey(params.receiverPublicKey, params.senderPublicKey)
//
// sign-client/src/controllers/engine.ts  authenticate()   [dapp side]
//   publicKey     = fresh x25519 keypair, stored under AUTH_PUBLIC_KEY_NAME
//   responseTopic = hashKey(publicKey)        -- hash of a PUBLIC key
//   request       = { authPayload: { type, chains, statement, aud, domain,
//                                    version, nonce, iat, exp, nbf, resources },
//                     requester: { publicKey, metadata }, expiryTimestamp }
//   sent on the pairing topic, sealed under the pairing symKey.
//
// sign-client/src/controllers/engine.ts  approveSessionAuthenticate()  [wallet]
//   receiverPublicKey = pendingRequest.requester.publicKey        (= gx here)
//   senderPublicKey   = fresh keypair                             (= gy here)
//   responseTopic     = hashKey(receiverPublicKey)
//   sessionTopic      = generateSharedKey(senderPublicKey, receiverPublicKey)
//   result            = { cacaos: auths,
//                         responder: { publicKey: senderPublicKey, metadata } }
//
//   NOTE the responder public key appears TWICE and the two are never compared:
//   once in the envelope (which decides the decryption key) and once inside the
//   ciphertext (which decides the session topic). They are modelled separately,
//   as gy_env and gy_pay.
//
// utils/src/cacao.ts
//   validateSignedCacao({ cacao }):
//     reconstructed = formatMessage(payload, payload.iss)   // payload IS cacao.p
//     return verifySignature(getDidAddress(payload.iss), reconstructed, ...)
//
//   The signed statement is therefore a deterministic function of the CACAO's
//   own payload, and the check is purely internal consistency: "the address
//   named in `iss` signed this text". It is NOT compared against the request.
//   Grepping the engine confirms it: `nonce` occurs only where a request is
//   built and in isValidAuthenticate's input validation, never in a comparison
//   against a returned CACAO, and neither `domain` nor `aud` is ever compared.
//   The model reflects that by signing the payload directly.
//
// ---------------------------------------------------------------------------
// THE BEHAVIOUR UNDER TEST
// ---------------------------------------------------------------------------
//
// engine.ts, dapp-side onAuthenticate:
//
//     const isValid = await validateSignedCacao({ cacao, projectId });
//     if (!isValid) {
//       this.client.logger.error(cacao, "Signature verification failed");
//       reject(getSdkError("SESSION_SETTLEMENT_FAILED", ...));
//     }
//     ... falls through ...
//     approvedChains  = [getNamespacedDidChainId(payload.iss)]
//     parsedAddress   = getDidAddress(payload.iss)
//     approvedMethods = getMethodsFromRecap(recap)          // from resources
//     const sessionTopic = await generateSharedKey(publicKey, responder.publicKey);
//     await this.client.core.relayer.subscribe(sessionTopic, { transportType });
//     await this.client.session.set(sessionTopic, session);
//
// reject() settles the promise but does not return, and onAuthenticate has no
// enclosing try/catch, so a CACAO that FAILS verification still reaches session
// construction. The approved account and methods are read from that same
// unvalidated payload. Store.set persists, so the session survives a restart.
// The wallet side of the identical check, ~190 lines below, does
// `await this.sendError(...)` and then `throw`.
//
// A check that does not halt is a SIGNVERIF without `?`. That is the whole
// difference between this model and 03-fix-verification.vp.
//
// ---------------------------------------------------------------------------
// PRECONDITION
// ---------------------------------------------------------------------------
//
// To address the response the attacker needs gx, which travels in the request
// under the pairing symKey. `leaks symkey` models the pairing URI being
// observed -- it is rendered as a QR code for the user to photograph, so it is
// exposed to anything that can see the screen. Over link mode this precondition
// disappears entirely: the request is a cleartext type 2 envelope in a URL.

attacker[active]

principal Wallet[
	knows private wsk            // the user's account key
	gw = PUBKEY(wsk)             // the did:pkh address, public knowledge
]

Wallet -> Dapp: [gw]

principal Dapp[
	knows private symkey         // pairing symKey, from the URI
	knows public c_ad
	knows public c_auth          // the wc_sessionAuthenticate method name
	knows public c_dom           // authPayload.domain and .aud
	generates x
	generates nonce              // authPayload.nonce
	generates txdata             // what the dapp sends once it believes it is authenticated
	gx = PUBKEY(x)               // requester.publicKey; responseTopic = sha256(gx)
	req = AEAD_ENC(symkey, CONCAT(c_auth, gx, nonce, c_dom), c_ad)
	leaks symkey
]

Dapp -> Wallet: req

principal Wallet[
	knows private symkey
	knows public c_ad
	knows public c_auth
	knows public c_dom
	knows public c_res           // authPayload.resources, carrying the recap
	generates y
	gy = PUBKEY(y)
	d_w = AEAD_DEC(symkey, req, c_ad)?
	tag_w, gx_w, nonce_w, dom_w = SPLIT(d_w)?
	_ = ASSERT(tag_w, c_auth)?
	// cacao.p: iss, domain, nonce, resources. formatMessage() is a function of
	// exactly this, so signing it directly is faithful.
	cp = CONCAT(gw, dom_w, nonce_w, c_res)
	cs = SIGN(wsk, cp)
	// The response is sealed under generateSharedKey(gy, gx) -- the session key.
	r = HKDF(nil, DH_KEX(gx_w, y), nil)
	// result = { cacaos: [{p, s}], responder: { publicKey } }
	resp = AEAD_ENC(r, CONCAT(cp, cs, gy), c_ad)
]

// Type 1 envelope: the senderPublicKey field travels in the clear.
Wallet -> Dapp: gy, resp

principal Dapp[
	// decode(): the decryption key comes from the envelope's senderPublicKey.
	r_d = HKDF(nil, DH_KEX(gy, x), nil)
	d_d = AEAD_DEC(r_d, resp, c_ad)?
	cp_d, cs_d, gy_pay = SPLIT(d_d)?
	// validateSignedCacao: reconstruct from the payload, verify against the
	// address the payload itself names. No comparison with nonce or c_dom.
	iss_d, dom_d, non_d, res_d = SPLIT(cp_d)?
	// The dapp is authenticating a PARTICULAR account -- it is re-authenticating
	// a known user, or matching the address against its own records. Without
	// this the queries below fail for an uninteresting reason: anyone may
	// answer an authenticate request with a CACAO for their OWN account, and
	// will of course hold the resulting session. That is connecting a wallet,
	// not impersonating one. This assertion is what makes the failure mean
	// "the dapp accepted a session for gw".
	_ = ASSERT(iss_d, gw)?
	// AS IMPLEMENTED: verification failure logs and rejects, then falls through.
	// No `?`, so the dapp does not halt on a bad signature.
	_ = SIGNVERIF(iss_d, cp_d, cs_d)
	// sessionTopic = generateSharedKey(publicKey, responder.publicKey) -- note
	// this uses the key from INSIDE the payload, not the envelope's.
	sess_d = HKDF(nil, DH_KEX(gy_pay, x), nil)
	sessreq = AEAD_ENC(sess_d, txdata, c_ad)
]

Dapp -> Wallet: sessreq

principal Wallet[
	_ = AEAD_DEC(r, sessreq, c_ad)?
]

queries[
	// Can the dapp be made to accept a response the wallet did not author?
	authentication? Wallet -> Dapp: resp
	// Does the attacker end up holding the session it establishes?
	confidentiality? txdata
]
\end{lstlisting}
%% --- END verifpal listing: 01-signature-bypass-vp-0 ---


\section{On soundness}
\noindent Verifpal is sound but incomplete. Every attack shown here is a genuine attack on the model as written; a query reported as holding means no attack was found within the search this run performed, which is never a proof that none exists.

\end{document}
